\section{Transport through Hodge-type integral models}
\label{sec:transport}

After taking a power of the original abelian variety, we transport a
rational Hodge class to an auxiliary variety by one quasi-isogeny on
special fibers, with the same equality in every
realization. The auxiliary variety can be chosen either of CM type or
with a prescribed admissible Hodge filtration. These choices serve
different parts of the proof, but share the tensor construction and the
rational normalization in Proposition~\ref{prop:transport} below.

Throughout this section, a polarization on a rational homology space is
written as a rational alternating form, using the fixed Betti
trivialization of the Tate line.  If its similitude character is $c$,
an alternating form $\xi:V^{2r}\to\Q$ of multiplier $c^r$ means that
\begin{equation}
 \xi(gv_1,\ldots,gv_{2r})=c(g)^r\xi(v_1,\ldots,v_{2r}).
 \label{eq:transport-multiplier}
\end{equation}
Such a form is regarded as a tensor with its codimension-$r$ Tate
twist.  All homomorphisms and polarizations may be rational: an
isomorphism in the isogeny category induces the usual isomorphisms on
cohomology, including its tensor constructions.

We say that filtered isocrystal data have \emph{finite local-field
descent} if they are obtained by scalar extension from an isocrystal
with its tensors over the maximal unramified subfield \(H_0\) of a
finite extension \(H/\Q_p\), together with a filtration over \(H\).
The scalar extension is to \(\Kzero H\). This condition will let us
apply weak admissibility over a finite local field.
For a group in a faithful representation, a tensor-preserving
identification means one respecting a finite family of rational tensors
whose stabilizer is that group, together with their comparison realizations.

\begin{proposition}[Transport]
\label{prop:transport}
Let $A/\Qbar$ be an abelian variety of positive dimension with good
reduction at $w$, and put $V=H_1(A(\C),\Q)$. Suppose
\[
 G\subset\GSp(V,\psi)
\]
is connected reductive, its Hodge homomorphism $h$ gives the
Hodge structure of $A$, and its conjugacy class $X$ makes this
inclusion a Hodge embedding. Let $\gamma$ be a rational Hodge
class of codimension $r$, represented on $V$ by an alternating
form satisfying \eqref{eq:transport-multiplier}.

Then there exist a positive integer $n$, an auxiliary
good-reduction abelian variety $B$, a quasi-isogeny
$a:A_0^n\to B_0$, and a rational Hodge class $\kappa$ of
codimension $r$ on $B$ such that, for the first-factor inclusion
$j:A_0\to A_0^n$,
\begin{equation}
 (aj)^*\kappa_0=\gamma_0.
 \label{eq:transport-class}
\end{equation}
The equality holds simultaneously in every $\ell$-adic
realization with $\ell\ne p$ and in the crystalline realization
after extension to $\Cw$. Either of the following two further
specifications can be imposed.
\begin{enumerate}
\item The variety $B$ is of CM type and is defined over $\Qbar$,
with its indicated good reduction at $w$.
\item Put $D_A=H^1_{\mathrm{cris}}(A_0/W(\F))[1/p]$. Suppose
$\mathcal F_A$ is a filtration in degrees $0,1$ on $D_A$, over a
finite extension of $\Kzero$, which is weakly admissible and has
finite local-field descent together with the specialized tensors
specifying the $G$-structure. Suppose
also that, after a geometric tensor-preserving identification, it
lies in the $G$-orbit of the original Hodge filtration. Then $B$
can be defined over a finite extension $K/\Kzero$ inside $\Cw$,
and $a$ can be chosen so that
\begin{equation}
 a^*\Fil^1H^1_{\mathrm{dR}}(B/K)
       =(\mathcal F_A^1)^{\oplus n}.
 \label{eq:transport-prescribed-filtration}
\end{equation}
Here the equality is made after a common finite extension and
crystalline--de Rham identification. To discuss Hodge classes on
this $B$, choose a complex embedding of an algebraic closure of
$K$ extending the fixed embedding of $\Qbar$.
\end{enumerate}
In the second case the realizations of $\kappa$ are supplied by
the canonical tensor sections of the Hodge-type model and the
polarization. In the first case these agree with the usual
realizations and specialization of its absolute Hodge class.
\end{proposition}

We prepare two ingredients for the proof. First, we enlarge the given
Hodge datum so that integral-model theorems apply at every prime while
its rational representation remains a sum of copies of \(V\).
Second, we realize a prescribed admissible filtration by choosing a
framed local Shimura point above its crystalline period and applying
integral uniformization.  The resulting abelian lift retains that
point's integral tensors and special-fiber marking.  After these
preparations, the proof constructs the transported tensor and normalizes
it using rational divisor intersections.

\subsection{An auxiliary strongly admissible datum}

We recall the hypotheses on integral models that we will use.
For a reductive group \(H_p/\Q_p\), a full Bruhat--Tits stabilizer
is the smooth group scheme fixing a point of its extended building.
Its identity component is the associated parahoric group scheme.
Now let \((H,X_H)\) be a Shimura datum over \(\Q\). In the terminology
of \cite[Definition~2.1.6]{KisinZhou2025}, the triple
\((H,X_H,\mathcal H)\) is \emph{strongly admissible} if \((H,X_H)\)
is of Hodge type, \(\mathcal H\) is a full stabilizer of \(H_{\Q_p}\)
with connected special fiber, the centralizer of a maximal
\(\Kzero\)-split torus of \(H_{\Kzero}\) is \(R\)-smooth, and the
relative root system over \(\Kzero\) is reduced when \(p=2\).  Here $R$-smoothness is the condition on tori used in
\cite[\S2.4]{KisinZhouEarlier}.  We will use it only for products and
extensions of restrictions of scalars of split tori, for which it
holds by \cite[Proposition~2.4.6]{KisinZhouEarlier}.

\begin{lemma}[Auxiliary Hodge datum]
\label{lem:auxiliary-datum}
Let
\[
 (G,X)\lhook\joinrel\longrightarrow
 (\GSp(V,\psi),X_{\mathrm{Siegel}})
\]
be a Hodge embedding, with $G$ connected and similitude character $c$.
For every prime $p$ there exist a totally real field $P$, a connected
reductive group $J$, and a Hodge embedding
\[
 (J,X_J)\lhook\joinrel\longrightarrow
 (\GSp(W,\psi_W),X'_{\mathrm{Siegel}})
\]
with the following properties.
\begin{enumerate}
\item There are inclusions of Shimura data induced by
\[
 G\longrightarrow J\longrightarrow\Res_{P/\Q}G_P,
\]
where the first map is diagonal and $J$ is the subgroup on which the
similitude multipliers in the restriction of scalars agree.
\item The group $J_{\Q_p}$ is quasi-split and has a full connected
stabilizer $\mathcal J$ making $(J,X_J,\mathcal J)$ strongly admissible.
\item As a rational $J$-representation,
\begin{equation}
 W\simeq (V\otimes_{\Q}P)^{\oplus m}
 \label{eq:transport-copies}
\end{equation}
for some positive integer $m$.  The form $\psi_W$ has multiplier the
common character $c$, and $\mathcal J\to\GL(W_{\Q_p})$ extends to a
closed immersion into $\GL(\Lambda_p)$ for a self-dual $\Z_p$-lattice
$\Lambda_p$.
\end{enumerate}
\end{lemma}

\begin{proof}
Choose a finite extension of $\Q_p$ splitting $G_{\Q_p}$.  Polynomial
approximation, together with Krasner's lemma
\cite[Propositions~7.60--7.61 and Corollary~7.62]{MilneANT}, gives a
totally real number field $P$ whose completions at $p$ split $G_{\Q_p}$.  For
example, approximate the minimal polynomial of a primitive element
of a splitting extension at $p$, and a polynomial with distinct real
roots at the real place.  The resulting number field can have just
one place above $p$; this additional condition is not needed.

Define the multiplier fiber product
\begin{equation}
 J=(\Res_{P/\Q}G_P)
       \mathbin{\times}_{\Res_{P/\Q}\mathbf G_m}\mathbf G_m,
 \label{eq:equal-multiplier-group}
\end{equation}
where the last map is diagonal.  The Hodge cocharacter, or its
inverse according to convention, pairs to $1$ with $c$.  Since $c$
factors through $G^{\mathrm{ab}}$, its character on this torus is
primitive.  Thus $\ker(G^{\mathrm{ab}}\to\mathbf G_m)$ is a torus,
and $\ker(c:G\to\mathbf G_m)$ is connected reductive.  The exact
sequence
\[
 1\longrightarrow\Res_{P/\Q}(\ker c)_P
 \longrightarrow J\xrightarrow{c}\mathbf G_m\longrightarrow1
\]
proves connectedness and reductivity of $J$.  Moreover,
\begin{equation}
 J^{\mathrm{der}}=\Res_{P/\Q}G_P^{\mathrm{der}}.
 \label{eq:auxiliary-derived}
\end{equation}

The action on $U=V\otimes_\Q P$, with alternating form
$\Tr_{P/\Q}\psi_P$, gives a symplectic embedding of $J$ with
similitude character $c$.  The diagonal Hodge homomorphism gives
$X_J$.  The adjoint Hodge-type and Cartan-involution axioms follow
componentwise from those for $X$.  We also check that extending the
totally real field introduces no rational adjoint factor on which
the Hodge homomorphism is trivial.  Write a rational simple adjoint
factor as $\Res_{L/\Q}H$, with $L$ totally real.  Every field factor
$L'$ of $P\otimes_\Q L$ is totally real, and every real embedding of
$L$ extends to a real embedding of $L'$.  Hence the existence of a
nontrivial Hodge component for $H$ persists in each resulting
rational factor.  Finally the diagonal weight is rational and the
trace form is a polarization.  This proves that $(J,X_J)$ is a
Shimura datum of Hodge type.

Here are the local conditions explicitly.  Put $P\otimes\Q_p=
\prod_{v\mid p}P_v$ and choose a split maximal torus $T_v$ of
$G_{P_v}$.  The character $c$ on $T_v$ has a cocharacter section:
a conjugate of the Hodge cocharacter pairs to $1$ with it, and all
cocharacters of a split torus are defined over $P_v$.  The same
argument applies to $G_{P_v}^{\mathrm{ab}}$.  Consequently the
maximal torus of $J_{\Q_p}$ obtained from the $T_v$, and the
abelianization of $J_{\Q_p}$, have the forms
\begin{align}
 T_J&\simeq\mathbf G_m\times
       \prod_{v\mid p}\Res_{P_v/\Q_p}\mathbf G_m^{\operatorname{rk}G-1},
       \label{eq:auxiliary-torus}\\
 J_{\Q_p}^{\mathrm{ab}}&\simeq\mathbf G_m\times
       \prod_{v\mid p}\Res_{P_v/\Q_p}
                    \mathbf G_m^{\dim G^{\mathrm{ab}}-1}.
       \label{eq:auxiliary-abelianization}
\end{align}
The split Borels over the $P_v$ give a Borel of $J_{\Q_p}$, so this
group is quasi-split.  Its derived group is a product of restrictions
of split groups by \eqref{eq:auxiliary-derived}.  The vertices from
the hyperspecial vertices in those split groups give a very special
parahoric of $J_{\Q_p}$.  The inertia coinvariants of the cocharacter
module in \eqref{eq:auxiliary-abelianization} are torsion-free, since
that module is a direct sum of permutation modules.  The full
stabilizer is therefore connected by
\cite[Lemma~4.2.4]{KisinZhouEarlier}.

The centralizer of a maximal $\Kzero$-split torus is obtained from
\eqref{eq:auxiliary-torus} by unramified base extension; its torus
factors are still induced split tori and $\mathbf G_m$.
It is $R$-smooth by
\cite[Proposition~2.4.6]{KisinZhouEarlier}.  Finally restrictions of
split root systems give reduced relative root systems, also at
$p=2$.  These are precisely the strong-admissibility checks in the
construction of \cite[Proposition~5.3.8]{KisinZhou2025}.  They use
the stipulated Hodge embedding, not the assertion that $G$ is the
full Mumford--Tate group of a particular point.

Apply \cite[Proposition~2.1.9]{KisinZhou2025}, starting with the
representation $U$ just constructed.  Its proof first tensors $U$
with an auxiliary totally real $\Q$-vector space, then takes direct
sums, and finally applies the symplectic self-duality construction.
The group acts trivially on that extra vector-space factor.  In the
last construction the dual representation twisted by $c$ is
identified with $U$ by its symplectic form.  Thus every resulting
rational summand is a copy of $U$, with the same multiplier.  This
gives \eqref{eq:transport-copies}, together with $\psi_W$ and
$\Lambda_p$.  The assertion of the cited proposition includes
$p=2$; its proof uses the reduced-relative-root condition through
\cite[Lemma~2.1.8]{KisinZhou2025}.
\end{proof}

Choose a sufficiently small neat prime-to-$p$ level $K^p$ and put
$K_p=\mathcal J(\Z_p)$.  The adapted embedding gives the integral
model $\mathcal S_K$ and its universal abelian scheme.  A rational
$\Z_{(p)}$-lattice can be chosen with completion $\Lambda_p$.
We use a finite family of tensors defining its stabilizer
$\mathcal J$, as in \cite[\S2.1.4]{KisinZhou2025}.
Any additional rational $J$-invariant tensors may be included in
this family after clearing denominators.  Indeed their integral
stabilizer condition already holds on the generic fiber; flatness
of $\mathcal J$ makes it hold on $\mathcal J$ itself.  Adding them
therefore changes neither the group scheme nor its tensor
trivializations.

The added tensors have compatible sections in every realization, as
follows.  Use the homological degree-one realizations in this paragraph,
that is, the duals of the corresponding $H^1$, so that their reference
representation is $W$.  For such a realization $H_\nu$ over its
coefficient field $k_\nu$, let
\[
 \mathcal P_\nu=\operatorname{Isom}_{\{s_\alpha\}}
                    (W\otimes k_\nu,H_\nu)
\]
be its torsor of frames preserving the original defining tensors.
For a rational $J$-invariant tensor $u$ and a frame $e$, set
$u_\nu=e(u)$.  Two frames differ by an element of $J$, so this
definition is independent of the frame and descends from a
trivialization of the torsor.  It also works integrally after clearing
the fixed denominator of $u$.

The Betti and \'etale frame torsors are supplied by the canonical
tensor sections of \cite[\S2.1.4]{KisinZhou2025}.  The construction
of the universal Breuil--Kisin--Fargues module and its structure
tensors gives the crystalline torsor and its compatibility with
those sections
\cite[\S\S2.1.13--2.1.14]{KisinZhou2025}.
Indeed, a comparison map carrying every $s_\alpha$ to its specified
realization carries a frame $e$ to the comparison of $e$; it therefore
carries $e(u)$ to the same tensor evaluated in the new frame.
Thus it preserves $u_\nu$ for every added invariant $u$, without
requiring $u$ to be an expression in the chosen defining tensors.
At a lift over $\Qbar$, these are the realizations of its rational
Hodge tensor, by Lemma~\ref{lem:comparison}.  Their crystalline
specializations are independent of the lift because they are
evaluations in the fixed crystalline frame torsor.

The same observation applies to isogenies.  One quasi-isogeny
preserving the defining tensors identifies their frame torsors in
every realization in which those tensors are preserved.  It
therefore preserves all the added invariant tensors as well.
This proves the tensor comparison and simultaneous preservation
that will be used below; flatness alone is used only to keep the
integral stabilizer unchanged.

Specification~\textup{(1)} of Proposition~\ref{prop:transport}
uses the special-point theorem
\cite[Theorem~2.2.7]{KisinZhou2025} directly.  The next lemma is
needed only for specification~\textup{(2)}, where the filtration
has been prescribed in advance.

\subsection{Realizing a prescribed admissible filtration}

For a point $x\in\mathcal S_K(\F)$ let $C_x$ be the corresponding
abelian variety and put
\[
 \mathbb D_x=H^1_{\mathrm{cris}}(C_x/W(\F)),\qquad
 D_x=\mathbb D_x[1/p].
\]
The crystalline structure tensors are denoted $s_{\alpha,0,x}$;
the tensors away from $p$ are denoted $s_{\alpha,\ell,x}$.
They are the specializations of the canonical tensor sections and
are independent of the chosen lift of $x$.

The \emph{tensor isogeny class} $\mathcal I_x$ consists of the points
$y\in\mathcal S_K(\F)$ for which a single quasi-isogeny
$C_x\to C_y$ preserves $s_{\alpha,0}$ and every
$s_{\alpha,\ell}$ for $\ell\ne p$.  We emphasize that this
definition concerns actual quasi-isogenies of abelian varieties.
It is the class defined in \cite[\S2.2.4]{KisinZhou2025}.

\begin{lemma}[Realization of an admissible filtration]
\label{lem:admissible-realization}
Let $(J,X_J,\mathcal J)$ be strongly admissible, with $J_{\Q_p}$
quasi-split, and choose the adapted embedding and neat level above.
Let $x\in\mathcal S_K(\F)$ have a lift to characteristic zero.
Suppose a filtration $\mathcal F$ on $D_x$, after a finite extension
of $\Kzero$, has the following properties:
\begin{enumerate}
\item the resulting filtered isocrystal is admissible;
\item geometrically its filtration is induced by a $J$-valued
cocharacter in the conjugacy class prescribed by $X_J$.
\end{enumerate}
Then, after extending that field finitely, there are $y\in\mathcal I_x$,
a tensor-preserving quasi-isogeny $b:C_x\to C_y$, and a lift
$\widetilde y$ over the extended valuation ring such that
\[
 b^*\Fil^1 H^1_{\mathrm{dR}}(C_{\widetilde y})=\mathcal F^1
\]
under the crystalline--de Rham identification.

In particular, condition \textup{(1)} follows if the filtered
isocrystal and its tensor data have finite local-field descent and
$\mathcal F$ is weakly admissible.
\end{lemma}

\begin{proof}
Write $k=\F$ and let $\breve E$ be the completed local reflex field.
We work over a common finite extension $K$ of $\Kzero$ containing
$\breve E$ and a field of definition of $\mathcal F$.  Every further
field extension below is finite.  Put $K_p=\mathcal J(\Z_p)$.

We first fix the relation between the group and cohomological
conventions.  Let $\rho:J_{\Q_p}\to\GL(W_{\Q_p})$ be the adapted
homological representation and let $\rho^\vee$ be its dual.  Choose
the tensor-preserving identification
$\mathbb D_x\simeq\Lambda_p^\vee\otimes_{\Z_p}W(k)$ of
\cite[\S2.1.14]{KisinZhou2025}.  The local shtuka parameter is
$b_x\in J(\Kzero)$, whereas crystalline Frobenius on this
cohomological Dieudonn\'e module is
\[
 F_x=\rho^\vee(\sigma(b_x))\circ\sigma.
\]
The additional $\sigma$ on $b_x$ comes from the comparison with
the vector shtuka: its special-fiber module is
$(\sigma^{-1})^*(\mathbb D_x^\vee)$
\cite[Example~2.3.4 and Remark~2.3.10]{PappasRapoport2024}.
Write $\{\mu\}$ for the class prescribed by $X_J$; then
$[b_x]\in B(J,\mu^{-1})$.

For precision, identify the general linear crystalline period
flag with the appropriate filtration variety on $D_x$ by its
marked Rapoport--Zink interpretation.  If a deformation
$\mathcal H$ has special-fiber marking
$q:\mathcal H_k\dashrightarrow C_x[p^\infty]$, this identification
sends its period to
$(q^*)^{-1}\Fil^1H^1_{\mathrm{dR}}(\mathcal H)$.
Pushout through the faithful Hodge representation gives the
commutative period square
\[
\begin{array}{ccc}
M_{\mathcal J,b_x,\mu,K_p}&\longrightarrow&M_{\GL,\rho(b_x),\rho(\mu)}\\
\pi_{\mathrm{GM}}\downarrow&&\downarrow\pi_{\mathrm{GM}}\\
\mathcal F_{J,\mu}&\longrightarrow&\operatorname{Flag}(D_x).
\end{array}
\]
The lower map is a closed immersion: the parabolic stabilizing
the faithful filtration intersects $J$ in the parabolic for
$\mu$.  Define $f_{\mathcal F}$ to be the unique $J$-flag with
image $\mathcal F$.  It exists by condition~\textup{(2)}.  The
square uses crystalline periods and the dual Frobenius descent
above, so no extension of $\sigma$ to the ramified filtration
field is chosen.  When $q=b^{-1}$ its lower right period is
exactly $b^*\Fil^1$.

\medskip\noindent
\emph{A framed point over the prescribed period.}
Admissibility of the faithful filtered isocrystal implies that
the image of $f_{\mathcal F}$ in the general linear period flag
variety is in its admissible locus.  To see this in the
cohomological coordinates of the square, let $U$ be a crystalline
representation realizing $(D_x,\mathcal F)$.  Its crystalline
comparison identifies the associated vector bundle away from
the untilt with $U\otimes\mathcal O$.  The filtered de Rham
comparison identifies the modification at the untilt with the
standard lattice $U\otimes B_{\mathrm{dR}}^+$.  They glue to the
trivial bundle $U\otimes\mathcal O$ on the Fargues--Fontaine
curve, which is the general linear admissibility condition.
Since $\rho$ is faithful,
$\rho^{-1}(Z_{\GL(W)})\subset Z_J$: an element acting as a scalar
commutes with the faithful image of $J$.  The admissible-locus
identity in \cite[Proposition~3.1.1(i), (3.1.9)]{PappasRapoport2024}
therefore puts $f_{\mathcal F}$ in
$\mathcal F_{J,\mu}^{\mathrm{adm}}$.

This locus is the image of the crystalline period map
\[
 \pi_{\mathrm{GM}}:
 M_{\mathcal J,b_x,\mu,K_p}
       \longrightarrow\mathcal F_{J,\mu}^{\mathrm{adm}}.
\]
The map is \'{e}tale
\cite[\S3.1.1, (3.1.4)]{PappasRapoport2024}.
Its fiber at the classical point $f_{\mathcal F}$ is nonempty;
an \'{e}tale rigid space over $K$ has a point over a finite
extension of $K$.  After making that extension, choose
\begin{equation}
 z\in M_{\mathcal J,b_x,\mu,K_p}(K),
 \qquad \pi_{\mathrm{GM}}(z)=f_{\mathcal F}.
 \label{eq:transport-framed-point}
\end{equation}
In particular $z$ supplies the integral level as well as the
rational isocrystal marking.  If $C$ is the completed algebraic
closure of $K$, its moduli description gives an actual
$\mathcal J$-torsor $\mathcal P_z$, a Frobenius isomorphism
bounded by $\mu$, and, for sufficiently large $a$, an isomorphism
\[
 i_a:\mathcal J_{Y_{[a,\infty)}(C^\flat)}
      \xrightarrow{\ \sim\ }
      \mathcal P_z|_{Y_{[a,\infty)}(C^\flat)},
 \qquad
 \varphi_{\mathcal P_z}\circ\varphi^*i_a
       =i_a\circ\varphi_{b_x}.
\]
These are precisely the torsor and annular framing in
\cite[\S3.1.1]{PappasRapoport2024}.  Its associated
pro-\'{e}tale $\mathcal J(\Z_p)$-torsor gives integral tensor
frames.  Thus the integral structure and Frobenius-compatible
framing are part of the chosen point, with all their tensor
compatibilities.

\medskip\noindent
\emph{Integral realization and its special fiber.}
Our adapted datum is of global Hodge type: the group scheme is
a full connected Bruhat--Tits stabilizer, and the chosen Hodge
embedding is integral on $\Lambda_p$.
The center also satisfies the split-rank hypothesis of
\cite[\S4.1]{PappasRapoport2024}.  Indeed its real subgroup of
multiplier one preserves the positive definite form
$\psi_W(v,h(i)w)$ and is compact.  Its real split rank is
therefore at most one, while the rational weight torus supplies
one split central direction over $\Q$; both split ranks are one.
The model here is the one in
\cite[\S2.1.11 and Theorem~2.1.12]{KisinZhou2025}, obtained
from this adapted Hodge embedding and satisfying the
Pappas--Rapoport characterization.  Condition $(U_x)$ holds
by \cite[Corollary~6.3]{GleasonLimXu2026}; see also
\cite[Remark~4.10.4]{PappasRapoport2024}.  Consequently
\cite[Theorem~4.10.6]{PappasRapoport2024} represents the integral
local Shimura variety by a normal formal scheme
$\mathfrak M_{\mathcal J,b_x,\mu}$ and gives the morphism
\[
 \Theta_x:\mathfrak M_{\mathcal J,b_x,\mu}
          \longrightarrow\widehat{\mathcal S_K}.
\]
More precisely, \cite[\S4.10.2, (4.10.2)--(4.10.6)]{PappasRapoport2024}
identifies this formal scheme with the tensor-compatible open and
closed part of the pullback of the Siegel Rapoport--Zink space.
This construction identifies the pullback of the universal
$\mathcal J$-shtuka with its tautological framed shtuka and the
faithful vector realization with the universal $p$-divisible group.

Its rigid generic fiber is the local Shimura variety
\cite[Proposition~3.5.1(i)]{PappasRapoport2024}.
The classical point $z$ of this rigid generic fiber gives
$\operatorname{Spf}\mathcal O_K\to\mathfrak M_{\mathcal J,b_x,\mu}$.
Indeed it lies in the generic fiber of a formal affine chart;
for a special formal chart one may take an affinoid member of
its standard generic-fiber covering containing $z$.
Evaluation of the formal chart's ring is integral, and its
ideal of definition maps to topologically nilpotent elements of
$\mathcal O_K$.  Let $z_0$ be the special point and put
$y=\Theta_x(z_0)$ and $\widetilde y=\Theta_x(z)$.
The formal map $\widetilde y$ gives a point of the integral model
over $\mathcal O_K$ by restricting to an affine neighborhood of
$y$.  Pullback of the universal abelian scheme gives
$C_{\widetilde y}$ and its $p$-divisible group $\mathcal H_z$.

We record explicitly the integral comparison at this point.
For a uniformizer $\pi$ of $K$, put
$\mathfrak S=W(k)[[u]]$, with $u\mapsto\pi$, and let
$E_\pi(u)$ generate its kernel.  The crystalline representation
and integral level of $z$ give the $\mathcal J$-Breuil--Kisin
torsor of \cite[\S3.5]{PappasRapoport2024}, with
\[
 \Phi:\varphi^*\mathcal P_{\mathrm{BK}}[1/E_\pi]
             \xrightarrow{\ \sim\ }
             \mathcal P_{\mathrm{BK}}[1/E_\pi].
\]
Its pushout to $\GL(\Lambda_p)$ is the Breuil--Kisin module of
$\mathcal H_z$.  Reduction at $u=0$, followed by the dual
Frobenius descent fixed above, therefore gives the actual
Dieudonn\'e lattice of $\mathcal H_{z,k}$.
Proposition~3.5.1 of \cite{PappasRapoport2024} compares the two
specializations of this already constructed framed point.  In
the coordinates of \cite[Theorem~2.2.5]{KisinZhou2025} it gives
\[
 g\in X_{\mathcal J}(\mu^{-1},b_x),\qquad
 \mathbb D(\mathcal H_{z,k})
       =\rho^\vee(\sigma(g))\mathbb D_x\subset D_x.
\]
Indeed the Frobenius matrix on this lattice, written in its
induced basis, is
\[
 \rho^\vee\!\left(
       \sigma(g^{-1}b_x\sigma(g))\right)\circ\sigma,
\]
which explains the same Frobenius descent in the
Deligne--Lusztig condition and in the lattice formula.

\medskip\noindent
\emph{The abelian marking and the filtration.}
Keep the identity prime-to-$p$ coordinate in the uniformization.
The description of its special fiber in
\cite[\S4.10.3]{PappasRapoport2024}, or equivalently the
isogeny-class map of \cite[Theorem~2.2.5]{KisinZhou2025}, gives
$y=i_x(g)$ and a single abelian quasi-isogeny
$b:C_x\to C_y$ preserving the crystalline tensors and every
prime-to-$p$ tensor.  Its crystalline marking is the one just
used for $\mathcal H_z$, because the uniformization is induced
by the marked Siegel Rapoport--Zink space.

The crystalline period map of that space is the Hodge filtration
of its universal deformation pulled back by its marking.
Compatibility of $\Theta_x$ with the tautological shtuka and
\eqref{eq:transport-framed-point} consequently give
\[
 b^*\Fil^1H^1_{\mathrm{dR}}(C_{\widetilde y}/K)
       =\mathcal F^1.
\]
The deformation criterion
\cite[Proposition~2.2.3]{KisinZhou2025} is satisfied concretely:
the filtration is of type $\mu$, and the crystalline structure
tensors compare with the integral \'{e}tale tensors in the
$\mathcal J(\Z_p)$-frames of $z$.  The bounded integral shtuka
and its annular framing were obtained before using this
criterion or comparing specializations.  All cited global
Hodge-type and strongly admissible results apply at $p=2$ as
well.

Finally, under finite local-field descent, the weak-admissibility
inequalities restrict to the descended Frobenius subobjects.
Colmez--Fontaine's theorem
\cite[Theorem~A]{ColmezFontaine2000} makes those descended data
admissible, and scalar extension supplies condition~\textup{(1)}.
\end{proof}

\subsection{Proof of the transport proposition}

\begin{proof}[Proof of Proposition~\ref{prop:transport}]
\medskip\noindent
\emph{The original point and its reduction.}
Apply Lemma~\ref{lem:auxiliary-datum}.  Choose a rational
$J$-equivariant identification
$W=(V\otimes_\Q P)^{\oplus m}$ and set $n=m[P:\Q]$.
Choose a basis of $P$ whose first vector is $1$.  The resulting
identification of the Hodge structure at the diagonal point $h$
with $V^{\oplus n}$ makes the first-factor map equal to
\begin{equation}
 v\longmapsto(v\otimes1,0,\ldots,0).
 \label{eq:transport-first-factor}
\end{equation}
The rational homology/isogeny correspondence
\cite[Corollary~6.9]{MilneShimura} therefore makes the abelian variety
\(C\) at that point isogenous to \(A^n\). Fix a rational isogeny
\(f:A^n\to C\) realizing this identification.

We explain the field and reduction assertions for this starting
point.  The Hodge embedding maps it to the Siegel point of the
polarized variety $C$ with finite level.  The abelian variety $C$
is defined over $\Qbar$: clearing a denominator in $f$ realizes
it as a quotient of $A^n$ by a finite subgroup of its torsion, and
every such subgroup is defined over $\Qbar$.  Its rational
polarization and its finite level also descend to $\Qbar$, since
homomorphisms between abelian varieties over $\Qbar$ do not
enlarge after extending algebraically closed fields
\cite[Corollary~20.4]{MilneAV}.  The Hodge
embedding induces a finite morphism onto its
image at sufficiently small finite level
\cite[Theorem~5.16]{MilneShimura}. Its fiber over a
$\Qbar$-point therefore consists of $\Qbar$-points.  This places
the diagonal point on the canonical model over $\Qbar$.

After a finite number-field extension the variety $C$ has good
reduction, since it is isogenous to $A^n$
\cite[Chapter~IV, Theorem~3.5]{MilneAVNotes}.  Its polarization and
prime-to-$p$ level extend after a further finite extension.
The integral Hodge-type model is the normalization of the closure
of the generic fiber in the adapted Siegel model.  The valuative
criterion for this finite normalization therefore extends our
point to its integral model.  Denote the reduction by $x$.
The rational isogeny $f$ specializes to a quasi-isogeny
$f_0:A_0^n\to C_x$.

\medskip\noindent
\emph{The trace tensor and its descent.}
Next extend the alternating form defining $\gamma$
$P$-multilinearly and take its trace.  Composing with projection
onto the first copy in \eqref{eq:transport-copies} gives
\begin{equation}
 \eta=\Tr_{P/\Q}(\gamma_P)\circ
              (\operatorname{pr}_1)^{\otimes 2r}
       \in\bigwedge^{2r}W^*.
 \label{eq:transport-trace-tensor}
\end{equation}
For $r=0$ this is interpreted as the scalar trace.
The equal-multiplier condition implies that $\eta$ has multiplier
$c^r$ under $J$.  On \eqref{eq:transport-first-factor}, its
restriction is
\begin{equation}
 j^*f^*\eta_C=[P:\Q]\gamma.
 \label{eq:transport-trace-restriction}
\end{equation}
The class $\eta_C$ is understood with its Tate twist.

We turn this semi-invariant into the untwisted tensors to which
the integral-model theorems directly apply.  Write
$\psi_W^{-1}\in\bigwedge^2W$ for the inverse symplectic
bivector, and include
\begin{equation}
 u_\psi=\psi_W\otimes\psi_W^{-1},\qquad
 u_\eta=\eta\otimes(\psi_W^{-1})^{\otimes r}
 \label{eq:transport-untwisted-tensors}
\end{equation}
among the structure tensors.  Both are $J$-invariant tensors in
ordinary tensor constructions of $W$ and $W^*$; no separately
trivialized Tate object is needed.  They can be made integral and
added to the defining family as explained above.  Their
crystalline and prime-to-$p$ realizations at $x$ are the
specializations from the original point, with the comparison
conventions fixed above.

\phantomsection
\label{par:transport-tensor-descent}
The entire finite family of structure tensors, including
$u_\psi$ and $u_\eta$, has finite local-field descent.  To see
this, choose a finite extension $L/\Q_p$ over whose valuation
ring the original number-field point is defined, and let $L_0$
be its maximal unramified subfield.  The canonical \'etale
tensor sections of \cite[\S2.1.4]{KisinZhou2025}, pulled back
to that point, are $\Gal(\overline L/L)$-invariant.  Applying
the tensor functor $D_{\mathrm{cris},L}$ to these morphisms
from the unit object places them in the crystalline tensor
spaces over $L_0$.  Their extension to $\Kzero$ is exactly
the family $s_{\alpha,0,x}$, by good-reduction comparison.
The map $f_0$ also descends after a finite residue-field
extension, so this descent is retained under its marking.
Taking a compositum with the finite field of definition of a
prescribed filtration therefore gives one finite local field
for all the data.  No tensor-preserving trivialization over
$L_0$ is being asserted or needed.

\medskip\noindent
\emph{The two choices of lift.}
For specification~(1), apply
\cite[Theorem~2.2.7]{KisinZhou2025} to $x$.  It gives a point
$y\in\mathcal I_x$ with a special-point lift and a
tensor-preserving quasi-isogeny $b:C_x\to C_y$.
The special-point lift supplies a CM abelian variety $B$.  Special
points on the canonical Shimura variety are algebraic over its
reflex field by the canonical-model theory
\cite[Definition~12.8 and Proposition~14.14]{MilneShimura}, so $B$
and its level data can be taken over $\Qbar$.  The given local lift specifies their embedding at $w$
and their good reduction $B_0=C_y$.

For specification~(2), identify $D_x$ with $D_A^{\oplus n}$ using
$f_0^*$ and prescribe the direct-sum filtration
$\mathcal F_A^{\oplus n}$.  It is weakly admissible, with the
required finite local-field descent.  It is also in the required
$J$-cocharacter orbit.  Indeed a geometric element of the
transported $G$ taking the original filtration on $D_A$ to
$\mathcal F_A$ acts diagonally on all the copies of $V$ in $W$;
its image in $J$ takes the original filtration on $D_x$ to the
prescribed one.  The tensor-preserving identifications used here
come from the initial point and comparison, so they include the
$J$-structure tensors as well.  Geometric conjugacy can be
realized over a finite extension, since the corresponding
isomorphism scheme is of finite type.
Lemma~\ref{lem:admissible-realization} now gives $y$, its lift
$B$, and a tensor-preserving quasi-isogeny $b:C_x\to B_0$ with
the prescribed marked filtration.  This proves
\eqref{eq:transport-prescribed-filtration} for
$a=bf_0$.

\medskip\noindent
\emph{A common rational normalization.}
It remains, in both specifications, to recover the twisted class
and normalize it by a single rational scalar.  Choose a complex
description of the lifted Shimura point.  Its rational marking by
$W$ supplies a rational Hodge class $\eta_B$ corresponding to
$\eta$, and a rational divisor class $\theta_B$
corresponding to $\psi_W$.  Write $\theta_C$ for the analogous
class on $C$ and put
\[
 \theta_*=f_0^*\theta_{C,0}\in\NS(A_0^n)\otimes_\Z\Q.
\]
The class $\theta_*$ is rational ample.  Tensor preservation of
$u_\psi$ says that, in each realization,
\begin{equation}
 a^*\theta_{B,0}=\lambda_v\theta_*,
 \label{eq:transport-local-similitudes}
\end{equation}
with a nonzero local scalar $\lambda_v$.  Let $D=n\dim A$.
Pairing with $\theta_*^{D-1}$ gives
\begin{equation}
 \lambda_v=
 \frac{\deg\bigl(a^*\theta_{B,0}\cdot\theta_*^{D-1}\bigr)}
      {\deg(\theta_*^D)}=:\lambda\in\Q^\times.
 \label{eq:transport-rational-similitude}
\end{equation}
Both intersection numbers refer to rational divisor classes on
the same abelian variety; their cycle-class traces are identical
in every realization.  The denominator is positive because
$\theta_*$ is ample, and the numerator is nonzero by
\eqref{eq:transport-local-similitudes}.  Thus all the scalars in
\eqref{eq:transport-local-similitudes}, including the crystalline
one, are the single rational number $\lambda$.

Preservation of $u_\eta$, together with
\eqref{eq:transport-rational-similitude}, then yields
\begin{equation}
 a^*\eta_{B,0}=\lambda^r f_0^*\eta_{C,0}
 \label{eq:transport-eta-scaling}
\end{equation}
simultaneously.  In fact the inverse bivector scales by
$\lambda^{-1}$, so that the equality for the untwisted tensor
in \eqref{eq:transport-untwisted-tensors} gives exactly
\eqref{eq:transport-eta-scaling}.  Set
\[
 \kappa=\frac{\lambda^{-r}}{[P:\Q]}\eta_B.
\]
Specializing \eqref{eq:transport-trace-restriction} and using
\eqref{eq:transport-eta-scaling} proves
\eqref{eq:transport-class}.

For the arbitrary local lift, the crystalline realization just
used is unambiguous: the canonical model supplies the untwisted
tensor $u_\eta$, and the algebraic polarization supplies its Tate
line and inverse bivector.  The comparison isomorphism and these
operations define $\eta_B$ in the crystalline realization and
agree with its canonical \'{e}tale tensors.  Consequently any
identity expressing this class by algebraic divisor operations
specializes with the same rational coefficients in all
realizations.  For the $\Qbar$-valued CM lift, these are also the
usual absolute-Hodge and de Rham realizations.  This completes
the two constructions.
\end{proof}

\begin{corollary}[CM transport of a Hodge class]
\label{cor:cm-transport}
Let \(A/\Qbar\) be an abelian variety of positive dimension with
good reduction at \(w\), and let
\(\gamma\in H^{2r}(A(\C),\Q(r))\) be a rational Hodge class,
where \(0\leq r\leq\dim A\). There are a positive integer \(n\),
a CM abelian variety \(B/\Qbar\) with good reduction at \(w\),
a rational Hodge class \(\kappa\) of codimension \(r\) on \(B\),
and a quasi-isogeny \(a:A_0^n\to B_0\) such that
\[
 \gamma_0=(aj)^*\kappa_0,
\]
where \(j:A_0\to A_0^n\) is the first-factor inclusion. This is an
equality of the usual absolute-Hodge specializations in every
prime-to-\(p\) realization and in crystalline cohomology extended
to \(\Cw\).
\end{corollary}

\begin{proof}
Choose a polarization \(\psi\) on \(V=H_1(A(\C),\Q)\), and let
\(G=\MT(V)\). The Mumford--Tate group is connected reductive and
is characterized by its Hodge tensors
\cite[Section~3, Propositions~3.4 and~3.6]{Deligne1982}.
The Hodge homomorphism and its \(G(\R)\)-conjugacy class define
the Mumford--Tate Shimura datum of Hodge type, with faithful
symplectic representation
\[
 G\hookrightarrow\GSp(V,\psi).
\]
The polarization supplies the rational weight and the Cartan
involution condition. The definition of the Mumford--Tate group
excludes an adjoint factor on which the Hodge homomorphism is
trivial. The class \(\gamma\), with its Tate twist, is fixed by
\(G\). Equivalently, its alternating form on \(V\) has multiplier
\(c^r\), where \(c\) is the similitude character. Thus the first
alternative of Proposition~\ref{prop:transport} applies and gives
the asserted data and simultaneous equality.
\end{proof}
